\begin{exercise}{104}
Let $p\colon\mathcal C\to\mathcal D$ be an inner fibration between
$\infty$-categories which induces a surjection on $0$-simplices and is a
Joyal equivalence. Show that $p$ is a trivial fibration.

\begin{proof}
We first prove that $p$ is an isofibration.
By Exercise~\ref{ex:89}, $hp\colon h\mathcal C\to h\mathcal D$ is an
equivalence of categories.
Let $x\in\mathcal C_0$ and let $e\colon p(x)\to d$ be an equivalence
in $\mathcal D$.
Surjectivity on vertices gives $x'\in\mathcal C_0$ with $p(x')=d$.
Fullness of $hp$ gives an edge $a\colon x\to x'$ with $[p(a)]=[e]$.
By Corollary~1.2.11 and Definition~1.2.3, there is a $2$-simplex
$\sigma$ in $\mathcal D$ with
\[
 d_2\sigma=p(a),\qquad d_1\sigma=e,\qquad d_0\sigma=s_0(d).
\]
The edges $a$ and $s_0(x')$ define a $\Lambda^2_1$-horn in
$\mathcal C$ lying over $\sigma$.
Since $p$ is an inner fibration, it has a filler $\widetilde\sigma$
over $\sigma$.
Its long edge $\widetilde e=d_1\widetilde\sigma\colon x\to x'$ satisfies
$p(\widetilde e)=e$.
As $hp$ reflects isomorphisms, $\widetilde e$ is an equivalence by
Definition~1.2.22.
This proves the equivalence-lifting condition of Definition~2.1.4.

Now $p$ is both an isofibration and a Joyal equivalence, so
Lemma~2.2.17 makes it a trivial fibration.
\end{proof}
\end{exercise}
